\subsection*{S18. Endpoint architecture and optimization study}

The study reuses the 1,280 public labels, 85 ARROW labels, 220 external identities
and original five-fold ARROW assignments. Source-surrogate weights, targets and
experimental labels are unchanged. It is retrospective, with adaptive development
before a final selection lock; it is not a new prospective validation. Final outer
predictions are made only after training-side selection is complete. The common
selection score averages public-source and ARROW-source validation MAE equally.
Three random inner splits each hold out 15\% of each eligible source, using seeds
$640601+j$ for $j=0,1,2$. Public training rows have loss weight 1 and ARROW rows
weight 3. No external label participates in these fits or inner decisions.

The initial 52-candidate list covers residual MLPs, dual-branch networks,
ridge-plus-neural residuals, pretrained-encoding heads and directed message passing.
Tabular residual blocks apply layer normalization, two linear transforms with GELU
activation and dropout, then add an identity skip. The dual architecture keeps a
separate 32-unit response branch before concatenating it with the structural branch.
The ridge model fits a training-only linear baseline with regularization 10 to the
scaled RDKit and response columns, then trains a residual network on the remaining
target. This is not a ridge model on seven selected responses with a tree residual.
The graph variants use Chemprop bond message passing with either summed and mean
atom embeddings plus atom count, or summed learned atomic outputs plus a global
correction. These readouts carry size information but do not impose exact additivity
of the full free-energy predictor. Frozen-encoder neural heads use separate
structure/encoding and response branches.

Candidate configurations vary widths 32--256, depth, dropout 0--0.2, learning rates
from $3\times10^{-4}$ to $3\times10^{-3}$, weight decay, squared or Huber loss,
and plateau-reduced or cosine schedules. AdamW and batch size 128 are used.
Ordinary neural fits have a 600-epoch cap. Median imputation, constant-column
removal and numerical scaling are fitted inside each training split; response
columns are retained even if constant. Binary Morgan bits remain unscaled.
The skewed Ipc descriptor is transformed by $\log(1+\max(\mathrm{Ipc},0))$ or
removed; BertzCT and Kappa1--3 use signed-log transforms. Physical counts, mass
and surface descriptors are not clipped to the training range. Count-based
fingerprints and later raw-count, log-count and no-fingerprint variants test whether
loss of multiplicity in binary Morgan features limits size response.

Each candidate is screened on split 0 at seed 11. The best two per family proceed
to three splits and seeds 11/29. The family winner is challenged with a 1,200-epoch
cosine schedule on the same three splits at seed 47, against its ordinary schedule.
The longer schedule is retained only if mean validation MAE decreases by more than
0.005 kcal mol$^{-1}$ and at least two of three pairs improve. A retained long
schedule is completed at seeds 11/29, giving nine inner fits. Thus this challenge
changes both duration and learning-rate trajectory, not duration alone. Fresh
outer fits use the median selected epoch and the inner-derived schedule. A width
boundary challenge expands selected 256-unit models to 512 units when a family
selected the upper width in at least four of six partitions; graph message dimension
is expanded as well. The same paired improvement rule determines retention.

TabM adds 12 configurations spanning three fingerprint treatments and four width,
depth, dropout and numerical-embedding settings. Each model uses 32 parameter-sharing
ensemble members; their mean is the prediction.\cite{tabm2025} The study uses TabM
0.0.3 and rtdl-num-embeddings 0.0.12. Representation refinements and the width
extension are retained as explicit selection variants, not relabelled as separate
architectures. Nine of 24 width challenges passed to nine-repeat completion.
All 81 retained wider-fit optima occur before 95\% of their training horizon.
Longer schedules were selected in all six corrected-MoLFormer partitions, five
dual and TabM partitions, three graph partitions, and two residual-MLP and ridge
partitions. Supplementary Fig.~12 shows the paired duration and width diagnostics.

The search contains 1,812 inner evaluation records, including 174 TabPFN context
evaluations. Gradient-based runs total 1,018,212 epochs; TabPFN's 2,112 internal
ensemble passes are not gradient epochs. Summed fitting time is 45.06 hours across
overlapping workers, not elapsed time or GPU-hours. The defined search and boundary
checks completed, but do not exhaust all architectures, regularization choices or
training-data interventions. Full candidate lists and per-partition decisions are
in Supplementary Data~6. Final comparisons retain every family, not only the global
selector. Matched no-response models inherit the response model's recipe, so the
comparison does not estimate each structural-only architecture's best achievable error.

\subsection*{S19. Pretrained representations and TabPFN configuration}

The earlier MoLFormer cache used evaluation mode without enabling the model's
\texttt{deterministic\_eval} option. Stochastic attention projections remained
active. All 1,585 study representations and 56 stress-series representations were
therefore regenerated from the same checkpoint with both settings enabled, and
all valid MoLFormer endpoint comparisons were refitted. Identical repeated encoding
gave zero change; changes in query order and batching produced representation
differences below $3\times10^{-5}$. Fresh encoding through endpoint prediction
agreed within $4.9\times10^{-6}$ kcal mol$^{-1}$. Corrected fixed-tree results
replace the earlier cache-based values in Fig.~5 and Supplementary Tables~11--12.
Legacy stochastic-cache branches remain labelled in the released diagnostic data
and are excluded from valid-model selection and the manuscript's performance
comparisons. Uni-Mol's conformer pipeline and the packaged ExtraTrees endpoint are
unaffected. Original encoder pretraining membership remains unaudited.

TabPFN-3.5 uses \texttt{tabpfn} 9.1.0 and full/Fast checkpoints from revision
\nolinkurl{06bf2ba35c80a92a3b9abb436b99cf49e7a0365e} of
\nolinkurl{Prior-Labs/tabpfn_3_5}.\cite{tabpfn35} Checkpoint SHA-256 values are
recorded in Supplementary Data~6. This specific software release is distinguished
from the earlier tabular foundation-model paper.\cite{tabpfn2025} We use its
pretrained weights without fine-tuning. Inputs retain training-fitted native
preprocessing, including feature subsampling/ranking and target transformations.
The use of a transformer endpoint therefore does not mean that all preprocessing
or the six source-surrogate families are neural.

Eight candidates combine full/Fast weights with four recipes: RDKit plus binary
Morgan bits; RDKit plus raw Morgan counts; RDKit alone; or RDKit plus the corrected
MoLFormer representation. Each has the 15 responses in its full condition.
The best two screen candidates receive repeated inner evaluation. Training-side
challenges compare one versus three copies of each eligible ARROW context row,
and eight versus 32 ensemble passes, using the same greater-than-0.005 and
two-of-three improvement rule. Context repetition occurs after splitting and
does not create new independent observations or equal weighted gradient fitting.
All selected recipes use the full checkpoint and three ARROW copies. The selected
ARROW folds 0--4 use respectively raw counts/8 passes, no fingerprint/32,
no fingerprint/8, raw counts/32, and MoLFormer without fingerprint/8. The external
fit uses MoLFormer without fingerprint and eight passes. These are a per-partition
selection policy, not one identical input recipe across folds.

The comprehensive policy selects the eligible option with lowest repeated inner
score independently in each partition. It selects TabPFN in all six partitions.
Five-seed final predictions use seeds 11, 29, 47, 71 and 101; controls remove only
the responses. Common-three-seed means address the unequal ensemble size relative
to the three-seed ExtraTrees fits. Every fitted-state replay and final identity/label
audit passed. Licensed TabPFN weights and fitted contexts are retained privately;
contexts embed training labels and are not public model-weight substitutes.

\subsection*{S20. Size holdout, homologue series and label influence}

The size split ranks molecules by descending heavy-atom count within public and
ARROW sources, breaking ties by identity. The largest 15\% comprise 192 public
and 13 ARROW molecules, leaving 1,160 training rows. Training heavy-atom counts
range from 0 to 17 and holdout counts from 8 to 37; 188 holdout molecules exceed
the global training maximum. Models use their locked external-fit recipes and
three fresh seeds. ExtraTrees is refitted on these 1,160 rows rather than evaluated
using release weights that have seen the held-out labels. No size score affects
selection. Since the held-out labels participated in earlier random-fold
development, intervals describe this exploratory split, not prospective selection-free
performance on new molecules. Supplementary Table~15 gives all valid family results.
The 13 ARROW molecules retain their original exclusion from supervised source
tables. No analogous source-disjoint guarantee is made for the 192 public molecules:
source models and source-OOF features remain fixed, with no size-specific source
refitting or exposure audit. Accordingly, this is a conditional endpoint comparison,
not a validation of fully source-unseen size extrapolation.

The unlabelled grid contains linear C1--C30 alkanes and Ac-(Gly)$_n$-NHMe and
Ac-(Ala)$_n$-NHMe for $n=0,\ldots,12$. The two zero-residue rows represent the
same N-methylacetamide molecule, so 56 rows are not 56 unique structures. All
positive-length peptides are absent from endpoint training; N-methylacetamide is
present. The longest glycine and alanine examples contain 53 and 65 heavy atoms,
respectively, beyond the maximum 37 in the original training pool. Binary Morgan
fingerprints are identical across $n=2,\ldots,12$ within either peptide family and
across C7--C30 alkanes. Count fingerprints, RDKit descriptors and some predicted
responses still change, so the entire input does not freeze. An initial neural
failure involved very large Ipc excursions; the present transformed-input study
avoids that numerical failure without establishing physical accuracy.

For ExtraTrees, each leaf value is an average of eligible training labels and the
ensemble output is an average of leaf values. Averaging each tree's minimum and
maximum leaf values gives an outer prediction bound, not necessarily a jointly
attainable extremum. The released ensemble bound is $[-32.5945,5.7565]$
kcal mol$^{-1}$. TabPFN's mean is a convex combination of finitely many bin means
on its fitted raw-target scale; averaging seedwise extrema gives the wider outer
bound $[-658.4946,649.9221]$. Its edge distributions can have unbounded tails
without an unbounded mean. Changing this bound is not a proof of size extensivity.

Only C17, C29 and C30 are absent from endpoint training. C14 and C16 labels of
2.16 and 1.22 kcal mol$^{-1}$ are faithful rounded means of SoluteML's linked
CompSol records (C14: 68956/68957; C16: 67229/67230). The revised CompSol archive
retains these binary-mixture records. Critical IUPAC--NIST evaluations identify
disagreement in the underlying solubility measurements and do not recommend a
best solubility value (\url{https://srdata.nist.gov/solubility/sol_detail.aspx?sysID=38_354}
and \url{https://srdata.nist.gov/solubility/sol_detail.aspx?sysID=38_389}).
We did not convert an extrapolated trend into a replacement label. A same-recipe
three-seed refit omitting only C14/C16 raises their predictions from 2.858/2.206
to 4.118/4.204 kcal mol$^{-1}$. This isolates substantial label influence but
does not establish that the refitted predictions are correct. All principal
architecture comparisons retain the original labels.

\subsection*{S21. Native predictive intervals}

For the five final external TabPFN fits, each predictive distribution is expressed
on the same raw kcal mol$^{-1}$ grid. Equal-weight mixture cumulative probabilities
are inverted to obtain central 80\% and 95\% intervals, rather than averaging
individual quantiles. Reconstructed means agree with stored predictions within
$1.4\times10^{-5}$ kcal mol$^{-1}$. The external 220 labels and strict 97 subset
are used to assess, not calibrate, coverage. At 80\%, 128/220 and 54/97 labels
are covered; at 95\%, 176/220 and 74/97 are covered. Supplementary Table~16 also
reports interval widths. These undercoverage rates preclude describing the native
intervals as calibrated reliability estimates. Seed dispersion and bootstrap
intervals for mean errors are different quantities. No peptide coverage can be
measured without reference free energies.
